\documentclass[11pt]{article}
\usepackage[review]{acl}
\usepackage{times}
\usepackage{latexsym}
\usepackage[T1]{fontenc}
\usepackage[utf8]{inputenc}
\usepackage{microtype}
\usepackage{inconsolata}
\usepackage{amsmath}
\usepackage{booktabs}
\usepackage{supertabular}

\title{Which Dictionary Should a Multilingual Lemmatizer Trust?\\
Source Priority and Rule Fallback in Multilingual Lemmatization}
\author{Anonymous submission}

\begin{document}
\maketitle

\begin{abstract}
Morphological dictionaries sometimes assign different lemmas to the same word.
We study which dictionary should take precedence and what happens when lookup
fails. For overlapping dictionaries we give an exact algorithm for choosing
an ordered subset on labeled training tokens. In a multilingual research
configuration, 15 of 98 languages with training evidence passed a separate
development check for a change to the declared source policy. We also audit
a 172-language catalog in which each packaged dictionary has one declared
resource. After screening Universal Dependencies test files against those
inputs, 121 languages retain 4,214,133 scorable tokens. Case-sensitive
accuracy averages 57.65\% across languages for copying, 64.42\% with the
packaged defaults, and 68.58\% with suffix rules disabled. The rule effect
varies with the training resource and language. Whether changing dictionary
priority improves the full system on independently collected documents of
the same language remains an open question.

\end{abstract}

\section{Introduction}

The Finnish word \emph{kirjakauppa} means \emph{bookshop}. A dictionary can
give \emph{kirjakauppa} as its headword, while the Finnish Turku Dependency
Treebank writes the lemma as \emph{kirja\#kauppa}, marking the compound
boundary. For exact-match evaluation, the boundary is part of the answer.
Finding an entry for a word therefore solves only half the problem. A
lemmatizer that draws on several resources must also decide whose analysis to
use. This decision recurs for capitalization, spelling variants and forms
shared by different lexemes. Two resources can both be internally consistent
and still disagree on the string a particular corpus expects.

External dictionaries are already useful for lemmatization.
\citet{bergmanis-goldwater-2019-data} use inflection tables to construct
training examples, and \citet{milintsevich-sirts-2021-enhancing} supply
analyses to a neural decoder at inference. We study a common implementation
choice with a different consequence for inference. Several sources are
compiled into a single table and the first source containing a form owns the
entry. Moving a source earlier changes every disputed form on which it takes
precedence over a different answer. Giving it lower priority preserves its
unique entries, which may still help on previously unseen words. Excluding it
also changes what happens on those unique entries. Accuracy measured
separately for each dictionary cannot tell us which action to take.

We study two decisions that occur at different stages of inference.
The first arises when two dictionaries cover a word. We formulate a
token-weighted objective for choosing an ordered subset and derive a subset
dynamic program that finds the best training policy. A separate development
comparison asks whether to change the declared priority. We report the outcome
for all 98 languages with training evidence, including candidates that failed
the check. In Turkish, ordering by each dictionary's standalone accuracy
would put the two leading sources in the opposite order from the optimum.

The second decision arises when lookup leaves a word unresolved. We isolate
the contribution of suffix rules in the project's 1.0.0 snapshot, whose 172
packaged dictionaries each name one input resource. This provides a control
in which source order is fixed by construction. It also exposes a problem
with catalog-wide evaluation. A resource labeled as a dictionary can itself
come from a UD test file. We trace its input before choosing which
test sentences count. Across the retained sample, turning on suffix rules
produces 148,380 corrections and 504,366 regressions relative to lookup
with the same dictionary. The language-level breakdown helps explain why
a single average would give a poor default for some resources.

These studies address different empirical quantities. The development
experiment measures decisions among multiple lexical sources. The packaged
comparison measures the effect of suffix fallback with the lexical content
held fixed. A further paired test is needed to estimate a full-system gain
from source reordering and its stability across corpora.

\section{Related work}

\paragraph{How lemmatizers use lexical evidence.}
LemmaGen induces ripple-down rules from form--lemma examples and uses them
to analyze new forms \citep{jursic-etal-2010-lemmagen}. Lemming jointly
infers morphological tags and lemmas \citep{muller-etal-2015-joint}.
Lematus evaluates the value of character-level sentence context on unseen and
ambiguous forms \citep{bergmanis-goldwater-2018-context}. An external resource
can enter such a model in at least two places. Paradigm entries can augment
training data when labeled sentences are scarce
\citep{bergmanis-goldwater-2019-data}. Alternatively, a decoder can learn
to copy characters from available candidate lemmas.
\citet{milintsevich-sirts-2021-enhancing} evaluate that strategy with
Apertium on 23 UD languages. Our predictor can decide at build time which
candidate will survive a conflict. It thus permits a direct comparison of
resource priorities with the lexical inputs held fixed.

\paragraph{Ordered decisions.}
First-match prediction has a long history. \citet{yarowsky-1994-decision}
ranks contextual evidence for lexical ambiguity resolution, and
\citet{angelino-etal-2018-rulelists} optimize regularized lists of rules for
categorical prediction. Our unit of selection is an entire dictionary. Its
condition is membership of the input form, and its answer varies with that
form. Once a prefix of dictionaries has claimed a set of forms, later
decisions can affect only the complement. We use this property to score all
orders and stopping points implicitly. Corpus frequencies, disagreement
between sources and the copy fallback determine which orders tie and whether
dropping a source can improve the objective.

\paragraph{Evaluation across languages and corpora.}
UD makes it possible to compare lemmatizers across treebanks with a common
representation \citep{nivre-etal-2020-universal}. Systems such as UDPipe and
Stanza also use sentence-level analyses
\citep{straka-strakova-2017-tokenizing,qi-etal-2020-stanza}. Their scores
quantify the cost or benefit of the whole modeling choice, including training
materials and context. \citet{toporkov-etal-2025-lemma} study how encoders and
in-context generation perform when supervised data is unavailable for the
target domain or language. Our transfer question concerns one learned quantity
within a fixed system, the priority assigned to sources. A fresh target corpus
and a fixed-order control would let us distinguish transferable lexical
evidence from adaptation to one treebank's editorial practice.

\section{Choosing an order of dictionaries}
\label{sec:method}

\subsection{Where a source can change a prediction}

The system takes pretokenized words and a language identifier. For an ordinary
token it tries corpus-derived form--lemma pairs, first by exact spelling and
then by language-aware lowercase spelling. An unresolved form can proceed
to the merged external vocabulary. If lookup fails, the system can apply a
learned suffix edit and ultimately return the original form. The available
stages depend on the language configuration. The catalog evaluation uses
ordinary tokens. The library also has paths for supplied contextual cues
and for entity mentions, which answer different questions.

The multi-source research configuration includes national dictionaries, analyses derived from
morphological tools, UniMorph paradigms \citep{kirov-etal-2018-unimorph}
and inflection tables extracted from Wiktionary
\citep{ylonen-2022-wiktextract}. Each source first yields a partial mapping
from an NFC-normalized surface form to one lemma. When an input resource
provides several analyses, the builder ranks them by the frequency of the
pair within that resource, followed by overall lemma frequency, shorter lemma
and source row order to break ties. This source-local decision is important
for reproducibility. Adding entries to another dictionary cannot break a tie
inside the first one.

In that configuration the declared source policy supplies a reference order before learning.
Its ordinary precedence starts with national dictionaries or Apertium,
followed by registered imports, UniMorph and Wiktionary, subject to declared
priorities. At build time the first source containing a form supplies the
value stored for that form. At query time the system looks up a single
compiled vocabulary. Consequently, the search over possible orders is an
offline cost. It does not require an extra model call or a traversal of the
source inventory for each input token.

\subsection{The training objective}

Let \(D_j(x)\) be the lemma that dictionary \(j\) gives form \(x\), and
\(C_j\) the set of forms it covers. Let \(c(x,y)\) count annotated training
tokens with form \(x\) and lemma \(y\). This count preserves distinct gold
lemmas for an ambiguous form and the frequency of each reading. For an ordered
subset \(\pi\), the prediction \(f_\pi(x)\) is the value of the first selected
dictionary covering \(x\), or \(x\) if no selected dictionary covers it.
We maximize the number of correctly lemmatized training tokens,
\begin{equation}
  Q(\pi) = \sum_x c\bigl(x,f_\pi(x)\bigr).
  \label{eq:objective}
\end{equation}
A dictionary that only reproduces forms already handled by copying can change
coverage without changing this objective. A source that rewrites them
incorrectly can reduce it.

For a set \(S\) of sources already placed, write
\(U(S)=\bigcup_{i\in S}C_i\). Appending source \(j\) changes the score
relative to copying on the previously unclaimed part of its domain.
\begin{equation}
  g(j,S)=\sum_{x\in C_j\setminus U(S)}
     \bigl[c(x,D_j(x))-c(x,x)\bigr].
  \label{eq:gain}
\end{equation}
The prefix order determines which answers were stored for forms in
\(U(S)\). Its membership alone determines which forms remain available to
\(j\). Hence the best cumulative gain \(F\) for each subset obeys
\begin{align}
  F(\emptyset)&=0, \nonumber\\
  F(S)&=\max_{j\in S}
     \{F(S\setminus\{j\})+g(j,S\setminus\{j\})\}.
  \label{eq:recurrence}
\end{align}
Removing the last dictionary of an optimal order leaves a prefix on the
smaller set. Replacing that prefix with a higher-scoring order would leave
the last gain unchanged and improve the total. This gives optimal
substructure. Taking the best state over \emph{all} subsets makes the
stopping point a property of the objective. The empty set represents
copying every word. There are \(2^n\) states and \(n2^{n-1}\) possible
append transitions for \(n\geq1\). The implementation precomputes gains by
coverage mask and limits exact search to 18 sources per language. At equal
correctness it prefers greater token coverage, greater form coverage, a
shorter list and then lexicographic source order.

This formulation explains why a large, individually accurate resource need
not go first. If a smaller dictionary covers disputed forms more accurately,
placing it first preserves the larger dictionary's unique entries. A
standalone score gives the larger resource credit for those unique entries
under either order. Equation~\ref{eq:gain} counts only the entries whose
predictions the particular decision can change.

\subsection{From a training optimum to a policy}

An optimizer may omit a dictionary whose placement makes no difference on
training data. Such an omission gives little reason to lose that source's
coverage of unseen words. For each omitted source we therefore examine the
best training score obtainable when that source must be included. An exclusion
is proposed only if this constraint costs at least 25 correct tokens or
0.01\% of scored training tokens, whichever is larger (rounding up). Sources
with weaker evidence remain available. The resulting candidate combines the
training-derived priorities and all strongly supported exclusions. If a
change is proposed, its complete specification is fixed before development.

The development comparison approximates the point at which an external
dictionary can influence a prediction. First it skips tokens covered by
training-derived exact or locale-aware lowercase corpus keys. On the remaining
tokens it tries the merged external table with exact and then lowercase
keys, followed by copying. This comparison uses neither suffix edits nor
entity processing. Accordingly, the optimizing objective, the development
comparison and the complete inference cascade have three distinct domains.

Let \(w\) count development tokens that a candidate corrects relative to
the declared policy, and \(l\) those it spoils. A change passes when
\(w-l \geq \max(5,\lceil N_{\rm dev}/10{,}000\rceil)\) and the exact
two-sided binomial form of McNemar's test on the \(w+l\) discordant tokens
gives \(p\leq0.05\). The effect threshold uses all scorable development
tokens even if few reach external lookup. The comparison accepts or rejects
the complete candidate in one step. We interpret the per-language test as
a decision rule, with correlated words and comparisons across languages in
mind when discussing the resulting count.

\subsection{How much can a priority change?}

Consider two policies run on the same \(N\) scorable tokens. Let \(D\) be
the positions at which they predict different lemma strings. At some of those
positions both strings may be wrong. Let \(w\) and \(l\) again count
corrections and regressions. For nonempty \(D\), the overall accuracy difference
can be written as
\begin{equation}
  \Delta_{\rm all}
  = \frac{|D|}{N}\,\frac{w-l}{|D|}.
  \label{eq:exposure}
\end{equation}
If the policies always predict the same lemma, their difference is zero.
Equation~\ref{eq:exposure} gives two ways for a promising lexical decision
to have a small effect on the full test. Few predictions may change, or the
changes may correct almost as many words as they spoil. Counting only
corrections discards the second possibility. Reporting accuracy only on
disputed words discards the first.

For pure reordering of fixed tables, dictionary membership is invariant.
An exact key claimed by a single source retains its answer in every order,
as does a key on which all sources agree. With earlier corpus lookups held
fixed, the potential effect is therefore confined to tokens that reach a
disputed key, including a lowercase key when exact lookup misses. Excluding
a source changes the interpretation. A word unique to that source can now
fall through to a suffix edit or to copying. This yields a concrete reason
to record both the stage reached and the prediction returned in a
controlled test of the complete system.

\section{Data and evaluation}
\label{sec:data}

\paragraph{Inputs to the priority experiment.}
The archived research configuration records one training optimization for
each of 98 languages and a subsequent development decision where a split
was available. It combines corpus lookups with external tables whose
provenance and order were fixed for that experiment. Some internal
development sets were formed earlier by partitioning an original UD test
split. We treat the resulting comparisons as development evidence. They can
explain why a policy passed an engineering decision rule, including why
the rule rejected apparently promising training candidates.

\paragraph{The 1.0.0 dictionaries.}
For the packaged-system experiment we verified the engine and the dictionary
artifacts from the same repository snapshot. The language roster has 172
entries, each with one declared upstream resource. For 101 entries that
resource is a UD treebank. Forty-three use a training file, 57 use a test
file, and one uses all three splits. The remaining 71 use resources such
as inflection paradigms and lexical databases. A lexical resource can contain
several files. We use the input list recorded for the individual package.
The source and runtime snapshot, checksums, language correspondence and
retained counts per treebank are recorded with the experiment.

\paragraph{Which test material remains?}
We start with the official UD 2.18 distribution \citep{ud218}. Of its
353 test treebanks, 330 have a file prefix corresponding to a packaged
language in the fixed code map. For every UD-derived dictionary we recover
the exact input files from its manifest and verify their digests. If an
upstream test file contributed to building the dictionary, we exclude
that file even if its revision changed. Within a treebank we check the
remaining document and sentence identifiers against the dictionary inputs.
Across treebanks of the same language we also match normalized word
sequences against the inputs and against already retained evaluation
sentences. Sequence comparison removes whitespace and folds case, then
normalizes to NFC. It includes punctuation and uses syntactic words in
their original order. This check can catch a repeated passage that was
renumbered when moving between corpora.

This procedure excludes 58 entire test files, 1,495 additional sentences
matching source text, and 7,265 repeated evaluation sentences. Of the 172
languages, 47 have only test files that supplied the dictionary. Three
further entries have no scorable words after exclusions. One language,
Turkmen, has no test file in the official distribution under its catalog
code. The resulting evaluation contains 261 treebanks, 121 languages and
4,214,133 words. The appendix retains a row for each of the original 172
languages. Counts for excluded languages belong to catalog coverage.

\paragraph{Predictions and uncertainty.}
We pass a sentence's syntactic words to the ordinary-token API in their
original case. We exclude withheld surface strings and missing gold lemmas.
A word with an available form and lemma contributes one case-sensitive
exact-match observation. Each word receives three predictions from the same
input. These are its surface form, the default package output, and the output of the
same engine and dictionary with suffix guessing disabled. User dictionaries
are disabled throughout. Gold lexical annotations and part-of-speech tags
are unavailable to the predictor. Tokens outside the scored subset stay
in their original sentence positions.

We report the pooled fraction of correct responses and the unweighted mean
of per-language accuracies, pooling retained treebanks within each language.
For a paired difference we resample blocks within each fixed treebank.
When a document boundary is available we use the document as the block.
Elsewhere we use the sentence. The 95\% interval uses 2,000 percentile
bootstrap replicates with seed 20261118 \citep{dror-etal-2018-hitchhikers}.
A resource-stratified analysis uses the same retained words and the same
three predictions.

\section{Results}

\subsection{When does research development evidence support a change?}
\label{sec:dev}

\begin{table}[t]
\centering\small
\begin{tabular*}{\columnwidth}{@{\extracolsep{\fill}}lr@{}}
\toprule
Outcome of policy selection & Languages \\
\midrule
Change confirmed on development data & 15 \\
Candidate evaluated, change declined & 38 \\
No change proposed from training data & 42 \\
Development split unavailable & 3 \\
\midrule
Languages with training evidence & 98 \\
\bottomrule
\end{tabular*}
\caption{Decisions about external dictionaries. Of the 15 confirmed policies,
six specify a new priority order and nine specify exclusions.}
\label{tab:decisions}
\end{table}

Table~\ref{tab:decisions} traces the 98 languages with annotated training
evidence. The procedure confirmed 15 changes to the declared policy. A
training optimum often supplied no actionable change, or its proposed change
failed on development data. The latter category includes candidates that
made accuracy worse and candidates for which the available sample could not
distinguish the two policies. A language without development data keeps its
declared policy under this procedure.

\paragraph{Why the best single dictionary goes second in Turkish.}
The training record for Turkish contains 166,521 annotated tokens and six
external sources. Used alone, with copying for missing entries, simplemma
scores 88.93\% and spaCy's lookup table scores 78.19\%. Yet the exact optimum
puts spaCy first and simplemma second, followed by Wiktionary and Zeyrek.
The fixed reference order yields 128,023 correct training tokens, compared
with 148,471 at the optimum. The two orders give different jobs to the
same resource. As the only dictionary, spaCy must cover a broad sample of
Turkish words. As the first dictionary in a cascade, it can supply its
preferred analysis on overlaps while simplemma still answers words it misses.

The two sources omitted at the training optimum also illustrate the
difference between an optimizer's tie-break and evidence for removing a
resource. The best order forced to include the national dictionary loses
four training tokens. Forcing UniMorph into an order costs zero. Neither
loss reaches the exclusion threshold of 25 tokens in this corpus. Both
sources therefore remain available behind the learned priority prefix in
the candidate tested on development data. Calling the entire optimized
training gain a benefit of dropping dictionaries would misdescribe the
intervention that was actually tested.

\begin{table}[t]
\centering\small
\begin{tabular*}{\columnwidth}{@{\extracolsep{\fill}}lrrrr@{}}
\toprule
Language & Dev & Reach & Wins & Losses \\
\midrule
Turkish & 6,995 & 872 & 69 & 38 \\
Spanish & 53,476 & 2,567 & 67 & 20 \\
English & 28,115 & 1,652 & 51 & 51 \\
Irish & 9,999 & 1,180 & 10 & 4 \\
Czech & 384,431 & 14,489 & 224 & 638 \\
\bottomrule
\end{tabular*}
\caption{Illustrative development decisions from the same confirmation
record. Reach counts tokens left after corpus lookup. Wins and losses compare
the training candidate with declared priorities. Turkish and Spanish passed.}
\label{tab:cases}
\end{table}

Turkish illustrates what a positive result amounts to
(Table~\ref{tab:cases}). Corpus lookup left 872 of 6,995 development tokens
to the lexical comparison. Moving sources within the order raised the number
correct on those tokens from 644 to 675 (\(p=0.0035\)). The 69 corrections
and 38 newly introduced errors show that each policy had words on which it
was preferable. The outcome establishes the direction of the paired
development difference for Turkish. It leaves open how often the same
changes would occur under full inference or recur on another collection.

Other languages give different answers to the same decision. The English
candidate exchanged 51 improvements for 51 regressions. In Irish, ten
improvements and four regressions cleared the minimum net count, but the
paired test gave \(p=0.180\). In Czech a candidate that included a
training-based exclusion produced 638 regressions and 224 improvements.
Its rejection prevented a substantial development loss. These outcomes
explain why the accepted subset alone gives a misleading impression of how
reliably a source policy can be learned.

\subsection{What do the packaged defaults contribute?}
\label{sec:catalog}

\begin{table}[t]
\centering\small
\begin{tabular*}{\columnwidth}{@{\extracolsep{\fill}}llrrr@{}}\toprule
Resource & Mean & Copy & Rules off & Full \\\midrule
All & Macro & 57.65 & 68.58 & 64.42 \\
 & Micro & 58.93 & 69.55 & 61.10 \\
\addlinespace[2pt]
UD-derived & Macro & 58.42 & 75.16 & 75.16 \\
 & Micro & 56.27 & 76.70 & 76.70 \\
\addlinespace[2pt]
External & Macro & 57.08 & 63.62 & 56.32 \\
 & Micro & 59.59 & 67.79 & 57.27 \\
\bottomrule\end{tabular*}

\caption{Case-sensitive token accuracy (\%). The same dictionary and test
words underlie each row. The UD-derived stratum has 52 scored languages
and the external-resource stratum 69.}
\label{tab:accuracy}
\end{table}

Table~\ref{tab:accuracy} reports the packaged-system comparison. The
default gains 6.76 macro accuracy points over copying (paired 95\% interval
6.61 to 6.89). The difference on pooled tokens is 2.17 points. These
aggregates weight the catalog differently. Relative to copying, 85
languages gain correct answers, 33 lose them and three tie. Across all
tokens the system corrects 715,637 copy errors and spoils 624,391 correct
copies. Looking up a form can correct a large number of inflections while
also changing a substantial number of words that already have their lemma
as the surface form.

Disabling suffix rules yields 68.58\% macro accuracy and 69.55\% pooled
accuracy. Enabling them changes the macro score by $-4.17$ points
(paired 95\% interval $-4.27$ to $-4.07$). The corresponding pooled
difference is $-8.45$ points. In eight languages the rules improve the
number of correct predictions, in 47 they reduce it, and in 66 it stays
the same. These counts describe accuracy ties. Two systems can have equal
numbers correct while producing different strings on individual tokens.

\paragraph{Where does the rule effect arise?}
All 52 scored languages with a UD-derived dictionary have the same number
correct with and without suffix rules. Their macro accuracy is 75.16\%,
compared with 58.42\% for copying. Among the 69 languages with an external
resource, the defaults average 56.32\%, copying 57.08\%, and lookup with
rules disabled 63.62\%. The difference between rule settings in this group
is $-7.31$ macro points. Resource assignment and evaluation coverage vary
with the language. The gap between the two groups therefore cannot by
itself identify the effect of replacing one source type with the other.

The paired counts show why a single choice of suffix settings has uneven
consequences. On 195,423 English words the rules correct 1,855 lookup
errors and introduce 197 new errors. In Spanish they correct 887 and
introduce 6,230 on 96,925 words. In Turkish there are
8,608 corrections and 10,851 regressions on 99,232 words. These languages
all have paradigms from UniMorph in their packaged dictionary. A shared
upstream resource family leaves considerable variation in what its suffix
edits do on running text.

\paragraph{Implications for a source comparison.}
Changing a source can alter the set of words that receives a suffix
prediction. A weaker lexical candidate might even improve accuracy if it
prevents a particularly harmful fallback. Conversely, excluding a source can
sacrifice entries on which its suffix replacement performs worse. A
full-system test of priority needs both policies evaluated with the same
fallback, plus a record of the positions where policy changes reach it.
The current packaged comparison locates the importance of that control.
No package contains the multiple source choices needed to estimate the
priority effect directly.

\subsection{What counts as a wrong lemma?}

The Finnish compound in the introduction illustrates an error category
that a better priority order alone may leave unresolved. If every source
gives the unsegmented headword, reordering can never produce the treebank's
boundary symbol. The archived development audit identified otherwise
matching lemmas that differed only in compound markers. These accounted
for 3.93 percentage points of all scorable development tokens in Finnish
TDT, 2.18 in Estonian EDT and 1.72 in Dutch LassySmall. In Finnish, 574
of the marker-only errors came from external lookup, 124 from the suffix
rule and 19 from copying. The error occurs at multiple stages, which
makes its linguistic description more informative than classifying every
instance simply as a dictionary miss.

The audit tested a separate repair. From training lemmas with marked
components it inferred the treebank's boundary marker and annotated
external compound lemmas with an unambiguous decomposition. On Finnish
development data this change corrected 400 tokens and introduced 21
regressions. These counts belong to the compound adapter's development
comparison. They measure a change to the contents of lexical entries.
The source-order comparison measures a choice among available entries.
A matched experiment must freeze this adapter along with the source
tables, especially across corpora that encode compounds differently.

\section{Does a learned policy travel?}
\label{sec:transfer}

The development decisions give us a concrete transfer question.
\emph{With lexical resources fixed, does an order learned on one corpus
improve token lemmatization on an independently collected corpus of the same
language?} The distinction between a change of corpus and a change of source
order is essential. Measuring one frozen system on two treebanks mixes the
difficulty of their vocabularies and the conventions of their annotators.
For example, the Korean GSD and KAIST treebanks mark morpheme boundaries in
lemmas with a plus sign, while Korean PUD omits those boundaries. A shift
in agreement with these conventions can register as a substantial loss
under exact match even when the lexical analysis is useful.

\paragraph{A paired in-domain comparison.}
The first contrast would run the declared order and the development-confirmed
policy through the same complete system on the same test tokens.
The main outcome would be the mean within-language accuracy difference
across all languages with training evidence and scorable evaluation data,
including those whose policy reverted to the declared order. A complementary
analysis would compare the best single dictionary, a greedy order chosen by
marginal training gain, the exact training candidate and the candidate after
development confirmation. Fitting these alternatives would use training data
only, reserving development data for the predefined confirmation step.
Holding the corpus tables, suffix rules, resource versions and enabled
stages fixed would isolate the source decision. A comparison with a
competitor supplying lexical material could also remove that material from
the cascade to expose the dependence.

\paragraph{A paired target-corpus comparison.}
For each eligible language we would freeze both policies and all other
components before inspecting target labels. We would then apply both to
the same independently sourced documents. Target eligibility would depend
on available lemma annotations, comparable token boundaries and verifiable
document provenance. Document identifiers and normalized sentence matches
would identify shared material with the source splits and earlier evaluation
sets. Where the target convention differs, the primary score would retain
the original labels. A prespecified diagnostic could score a documented
convention separately. Languages with no independent target or with too
few disputed tokens would remain visible in the coverage report.

\paragraph{Interpreting the outcome.}
Besides overall exact match, the comparison would report reachable tokens,
forms with disagreeing sources, forms whose gold lemma differs from their
spelling, and corrections and regressions. For an order-only change, both
orders agree on every exact key with a single proposed lemma. An exclusion
can also change the result on a form that occurs in just one source.
Separating these cases would distinguish a priority that resolves conflicts
from one that gains by suppressing harmful rewrites. Paired resampling by
document within language, using sentence blocks where document boundaries are
unavailable, would yield intervals that respect the pairing
\citep{dror-etal-2018-hitchhikers}. A positive in-domain effect alongside
a precisely estimated target effect near zero would support a narrower claim
about corpus adaptation. A wide interval on the target would instead leave
the direction and size of transfer unresolved.

\section{Conclusion}

Dictionary precedence can be learned exactly on a token-weighted training
objective. In the research configuration a separate development check
supported 15 policy changes across 98 languages with training evidence.
In the 1.0.0 package catalog, tracing dictionary inputs to their source
removed 58 test files from evaluation. On the remaining scorable words,
turning off suffix rules raised mean language accuracy from 64.42\% to
68.58\%. These findings establish how source overlap affects what can be
measured and how one fallback changes a prediction. The effect of a
multi-source policy on independent documents remains to be measured with
paired predictions from both frozen orders.

\section*{Limitations}

The 121 languages with usable test data represent 70\% of this package
catalog. Some of the smallest language resources were built from the only
available annotated test collection. Counting their own entries as test
gold would give a different scientific quantity. Eleven scored languages
contribute fewer than 1,000 words. Genre, annotation conventions, word
segmentation and the amount of masked text also differ. The bootstrap
conditions on the retained languages and treebanks and assumes independent
blocks within each treebank. For corpora without document boundaries,
repeated constructions can span several sentence blocks.

The provenance screen checks the inputs declared by each dictionary.
Reworded passages, translations and undocumented sharing upstream of a
lexical database may evade a string-based screen. Earlier engineering work
also used UD. The benchmark describes a fixed catalog on available text.
An assessment of transfer requires documents collected independently of
source-policy fitting and earlier error inspection.

The exact training optimum concerns NFC exact-form evidence and a copy
fallback. The development comparison additionally masks known corpus forms
and allows lowercase lexical lookup. Its per-language McNemar decisions
treat discordant tokens as units, despite correlation between repeated forms.
They provide no family-wise guarantee. The package ablation holds one
declared lexical resource fixed per language. A direct comparison of
multi-source policies on the current engine still requires a separately
frozen inventory and matched test predictions.

\section*{Ethical considerations}

The dictionaries have different licenses and attribution requirements.
For each package the evidence inventory records the upstream source and
license. The reproduction procedure verifies package bytes and obtains UD
from its distributor. It preserves the corpora's restrictions on withheld
text. The accompanying language and treebank ledgers contain counts and
digests. Their coverage record makes visible where available annotations
were already used to construct a dictionary and where a separate evaluation
would require further work, especially for languages with little labeled text.

\bibliography{references}

\appendix
\section{Language-level results}
\label{app:catalog}
The appendix enumerates all 172 dictionaries in snapshot 1.0.0. Sets gives
scored treebanks divided by in-scope treebanks. Words counts scored syntactic
tokens. Copy, Off and Full give exact-match accuracy (\%) for copying, the
packaged engine with suffix rules disabled and the packaged defaults.
\emph{Input file} means that every available test file had contributed to the
dictionary. \emph{No test} means that the official archive had no test file
for the code. \emph{No gold} denotes an entry with no scorable words after
exclusions. The electronic treebank ledger records file digests, paired
counts, overlap exclusions and missing annotations.

\begingroup\small
\setlength{\tabcolsep}{3pt}\renewcommand{\arraystretch}{0.95}
\newcommand{\cathead}{\toprule Language & Sets & Words & Copy & Off & Full \\\midrule}
\tablefirsthead{\cathead}\tablehead{\cathead}
\tabletail{\bottomrule}\tablelasttail{\bottomrule}
\begin{supertabular*}{\columnwidth}{@{\extracolsep{\fill}}p{.31\columnwidth}rrrrr@{}}
\raggedright (Modern) Greek {\scriptsize(ell)} & 8/8 & 42,038 & 44.7 & 50.4 & 43.6 \\
\raggedright Abaza {\scriptsize(abq)} & 0/1 & 0 & \multicolumn{3}{c}{Input file} \\
\raggedright Abkhaz {\scriptsize(abk)} & 0/1 & 0 & \multicolumn{3}{c}{Input file} \\
\raggedright Afrikaans {\scriptsize(afr)} & 1/1 & 10,063 & 78.9 & 95.2 & 95.9 \\
\raggedright Akkadian {\scriptsize(akk)} & 1/2 & 1,814 & 9.4 & 9.6 & 7.9 \\
\raggedright Akuntsu {\scriptsize(aqz)} & 0/1 & 0 & \multicolumn{3}{c}{Input file} \\
\raggedright Albanian {\scriptsize(sqi)} & 2/2 & 1,293 & 51.8 & 59.3 & 57.0 \\
\raggedright Amharic {\scriptsize(amh)} & 1/1 & 9,573 & 76.0 & 71.0 & 66.1 \\
\raggedright Ancient Greek {\scriptsize(grc)} & 3/3 & 44,503 & 23.3 & 24.0 & 22.6 \\
\raggedright Ancient Hebrew {\scriptsize(hbo)} & 1/1 & 14,972 & 13.0 & 13.1 & 14.0 \\
\raggedright Apurina {\scriptsize(apu)} & 0/1 & 0 & \multicolumn{3}{c}{Input file} \\
\raggedright Arabic {\scriptsize(arb)} & 1/3 & 28,249 & 13.2 & 65.5 & 63.5 \\
\raggedright Armenian {\scriptsize(hye)} & 2/2 & 20,819 & 52.5 & 61.5 & 71.0 \\
\raggedright Assyrian {\scriptsize(aii)} & 0/1 & 0 & \multicolumn{3}{c}{Input file} \\
\raggedright Azerbaijani {\scriptsize(aze)} & 1/1 & 873 & 49.9 & 54.4 & 56.8 \\
\raggedright Bambara {\scriptsize(bam)} & 0/1 & 0 & \multicolumn{3}{c}{Input file} \\
\raggedright Basque {\scriptsize(eus)} & 1/1 & 24,343 & 48.7 & 63.8 & 69.1 \\
\raggedright Belarusian {\scriptsize(bel)} & 1/1 & 15,958 & 49.0 & 49.6 & 40.7 \\
\raggedright Bengali {\scriptsize(ben)} & 1/1 & 320 & 65.6 & 73.4 & 51.2 \\
\raggedright Bhojpuri {\scriptsize(bho)} & 0/1 & 0 & \multicolumn{3}{c}{Input file} \\
\raggedright Bokota {\scriptsize(sab)} & 0/1 & 0 & \multicolumn{3}{c}{Input file} \\
\raggedright Bororo {\scriptsize(bor)} & 1/1 & 13,924 & 60.3 & 87.8 & 88.1 \\
\raggedright Brahui {\scriptsize(brh)} & 0/1 & 0 & \multicolumn{3}{c}{Input file} \\
\raggedright Breton {\scriptsize(bre)} & 1/1 & 9,996 & 56.9 & 54.0 & 31.5 \\
\raggedright Bulgarian {\scriptsize(bul)} & 1/1 & 15,724 & 52.9 & 59.2 & 52.6 \\
\raggedright Buryat {\scriptsize(bxr)} & 1/1 & 9,960 & 54.5 & 56.9 & 57.2 \\
\raggedright Cantonese {\scriptsize(yue)} & 0/1 & 0 & \multicolumn{3}{c}{Input file} \\
\raggedright Cappadocian {\scriptsize(cpg)} & 1/2 & 816 & 51.2 & 51.2 & 50.9 \\
\raggedright Catalan {\scriptsize(cat)} & 1/1 & 59,571 & 68.5 & 72.1 & 52.0 \\
\raggedright Cebuano {\scriptsize(ceb)} & 1/1 & 1,517 & 60.4 & 60.4 & 59.9 \\
\raggedright Chintang {\scriptsize(ctn)} & 1/1 & 14,199 & 48.9 & 54.7 & 54.9 \\
\raggedright Classical Armenian {\scriptsize(xcl)} & 1/1 & 11,185 & 61.2 & 66.1 & 63.2 \\
\raggedright Classical Chinese {\scriptsize(lzh)} & 2/2 & 25,797 & 98.1 & 99.7 & 99.7 \\
\raggedright Coptic {\scriptsize(cop)} & 2/2 & 21,409 & 64.6 & 80.7 & 79.5 \\
\raggedright Croatian {\scriptsize(hrv)} & 1/1 & 24,249 & 47.4 & 84.8 & 83.2 \\
\raggedright Czech {\scriptsize(ces)} & 7/7 & 375,892 & 53.4 & 68.9 & 65.9 \\
\raggedright Danish {\scriptsize(dan)} & 1/1 & 10,023 & 66.8 & 69.7 & 65.7 \\
\raggedright Dutch {\scriptsize(nld)} & 2/2 & 39,961 & 72.8 & 73.2 & 53.3 \\
\raggedright Egyptian {\scriptsize(egy)} & 1/1 & 11,444 & 52.9 & 86.3 & 85.0 \\
\raggedright English {\scriptsize(eng)} & 11/14 & 195,423 & 74.4 & 76.7 & 77.2 \\
\raggedright Erzya {\scriptsize(myv)} & 1/1 & 10,471 & 50.3 & 69.2 & 79.7 \\
\raggedright Esperanto {\scriptsize(epo)} & 1/2 & 175 & 61.1 & 73.1 & 86.9 \\
\raggedright Estonian {\scriptsize(est)} & 2/2 & 61,488 & 48.6 & 57.3 & 53.5 \\
\raggedright Faroese {\scriptsize(fao)} & 2/2 & 11,039 & 58.5 & 67.1 & 60.1 \\
\raggedright Finnish {\scriptsize(fin)} & 4/4 & 70,280 & 44.0 & 66.7 & 64.7 \\
\raggedright French {\scriptsize(fra)} & 9/9 & 118,564 & 61.6 & 60.4 & 30.0 \\
\raggedright Galician {\scriptsize(glg)} & 2/3 & 39,831 & 64.3 & 84.3 & 92.4 \\
\raggedright Georgian {\scriptsize(kat)} & 2/2 & 20,225 & 57.8 & 70.4 & 79.6 \\
\raggedright German {\scriptsize(deu)} & 4/4 & 409,328 & 66.2 & 74.3 & 51.0 \\
\raggedright Gheg {\scriptsize(aln)} & 0/1 & 0 & \multicolumn{3}{c}{Input file} \\
\raggedright Gothic {\scriptsize(got)} & 1/1 & 10,052 & 41.4 & 54.1 & 34.9 \\
\raggedright Guajajara {\scriptsize(gub)} & 0/1 & 0 & \multicolumn{3}{c}{Input file} \\
\raggedright Gujarati {\scriptsize(guj)} & 2/2 & 1,964 & 88.2 & 87.2 & 47.0 \\
\raggedright Gwichin {\scriptsize(gwi)} & 0/1 & 0 & \multicolumn{3}{c}{Input file} \\
\raggedright Haitian Creole {\scriptsize(hat)} & 2/2 & 10,599 & 85.0 & 98.3 & 95.8 \\
\raggedright Hausa {\scriptsize(hau)} & 3/4 & 36,810 & 64.8 & 74.1 & 74.7 \\
\raggedright Hebrew {\scriptsize(heb)} & 4/4 & 37,102 & 73.0 & 73.7 & 61.0 \\
\raggedright Highland Puebla Nahuatl {\scriptsize(azz)} & 0/1 & 0 & \multicolumn{3}{c}{Input file} \\
\raggedright Hindi {\scriptsize(hin)} & 2/2 & 59,242 & 72.3 & 76.3 & 38.0 \\
\raggedright Hittite {\scriptsize(hit)} & 0/1 & 0 & \multicolumn{3}{c}{Input file} \\
\raggedright Hungarian {\scriptsize(hun)} & 1/1 & 10,446 & 66.1 & 78.8 & 75.4 \\
\raggedright Icelandic {\scriptsize(isl)} & 4/4 & 178,358 & 53.4 & 55.0 & 46.2 \\
\raggedright Ika {\scriptsize(arh)} & 0/1 & 0 & \multicolumn{3}{c}{Input file} \\
\raggedright Indonesian {\scriptsize(ind)} & 3/3 & 41,182 & 72.8 & 83.7 & 83.4 \\
\raggedright Irish {\scriptsize(gle)} & 3/3 & 30,275 & 63.6 & 65.7 & 59.3 \\
\raggedright Italian {\scriptsize(ita)} & 11/11 & 126,472 & 64.3 & 64.9 & 39.4 \\
\raggedright Japanese {\scriptsize(jpn)} & 2/6 & 41,820 & 79.9 & 79.4 & 76.1 \\
\raggedright Javanese {\scriptsize(jav)} & 0/1 & 0 & \multicolumn{3}{c}{Input file} \\
\raggedright Kaapor {\scriptsize(urb)} & 0/1 & 0 & \multicolumn{3}{c}{Input file} \\
\raggedright Kadiwéu {\scriptsize(kbc)} & 1/1 & 836 & 49.9 & 49.9 & 46.1 \\
\raggedright Kangri {\scriptsize(xnr)} & 0/1 & 0 & \multicolumn{3}{c}{Input file} \\
\raggedright Karelian {\scriptsize(krl)} & 1/1 & 3,093 & 43.3 & 43.7 & 40.0 \\
\raggedright Karo {\scriptsize(arr)} & 0/1 & 0 & \multicolumn{3}{c}{Input file} \\
\raggedright Kazakh {\scriptsize(kaz)} & 1/1 & 10,007 & 50.1 & 50.4 & 50.5 \\
\raggedright Khoekhoe {\scriptsize(naq)} & 1/1 & 6,638 & 70.3 & 86.4 & 95.0 \\
\raggedright Kiche {\scriptsize(quc)} & 0/1 & 0 & \multicolumn{3}{c}{Input file} \\
\raggedright Komi Permyak {\scriptsize(koi)} & 1/1 & 898 & 57.9 & 63.5 & 64.0 \\
\raggedright Komi Zyrian {\scriptsize(kpv)} & 2/2 & 10,349 & 56.5 & 59.5 & 56.4 \\
\raggedright Korean {\scriptsize(kor)} & 5/5 & 74,470 & 31.7 & 59.6 & 70.0 \\
\raggedright Kyrgyz {\scriptsize(kir)} & 2/2 & 12,004 & 49.5 & 70.1 & 78.1 \\
\raggedright Latin {\scriptsize(lat)} & 6/6 & 102,238 & 43.9 & 46.9 & 33.0 \\
\raggedright Latvian {\scriptsize(lav)} & 2/2 & 37,440 & 47.5 & 59.9 & 56.7 \\
\raggedright Ligurian {\scriptsize(lij)} & 1/1 & 6,683 & 66.2 & 66.3 & 63.7 \\
\raggedright Lithuanian {\scriptsize(lit)} & 2/2 & 13,013 & 44.4 & 49.2 & 56.3 \\
\raggedright Livvi {\scriptsize(olo)} & 1/1 & 1,485 & 46.7 & 76.7 & 72.0 \\
\raggedright Low German {\scriptsize(nds)} & 1/1 & 10,966 & 61.7 & 61.1 & 38.2 \\
\raggedright Macedonian {\scriptsize(mkd)} & 1/1 & 1,343 & 65.9 & 74.4 & 61.3 \\
\raggedright Maghrebi Arabic French {\scriptsize(qaf)} & 1/1 & 2,162 & 21.6 & 45.9 & 45.7 \\
\raggedright Malayalam {\scriptsize(mal)} & 0/1 & 0 & \multicolumn{3}{c}{Input file} \\
\raggedright Mandarin Chinese {\scriptsize(cmn)} & 6/7 & 82,393 & 98.9 & 82.9 & 82.1 \\
\raggedright Manx {\scriptsize(glv)} & 1/1 & 10,252 & 72.9 & 93.2 & 93.3 \\
\raggedright Marathi {\scriptsize(mar)} & 2/2 & 12,703 & 57.0 & 81.7 & 85.9 \\
\raggedright Mbya Guarani {\scriptsize(gun)} & 1/1 & 1,318 & 60.1 & 62.1 & 62.1 \\
\raggedright Middle Armenian {\scriptsize(axm)} & 0/1 & 0 & \multicolumn{3}{c}{Input file} \\
\raggedright Middle French {\scriptsize(frm)} & 2/2 & 20,670 & 48.5 & 50.8 & 28.7 \\
\raggedright Moksha {\scriptsize(mdf)} & 1/1 & 4,305 & 49.8 & 53.6 & 56.1 \\
\raggedright Munduruku {\scriptsize(myu)} & 0/1 & 0 & \multicolumn{3}{c}{Input file} \\
\raggedright Naga {\scriptsize(nmf)} & 1/1 & 2,890 & 55.4 & 61.0 & 63.4 \\
\raggedright Naija {\scriptsize(pcm)} & 1/1 & 14,187 & 94.0 & 98.3 & 98.0 \\
\raggedright Nenets {\scriptsize(yrk)} & 0/1 & 0 & \multicolumn{3}{c}{Input file} \\
\raggedright Nepali {\scriptsize(nep)} & 0/1 & 0 & \multicolumn{3}{c}{Input file} \\
\raggedright Nheengatu {\scriptsize(yrl)} & 1/1 & 9,898 & 68.4 & 74.3 & 75.0 \\
\raggedright Northern Kurdish {\scriptsize(kmr)} & 1/1 & 9,963 & 53.0 & 60.9 & 68.2 \\
\raggedright Northern Sami {\scriptsize(sme)} & 1/1 & 9,971 & 47.5 & 60.2 & 53.6 \\
\raggedright Northwest Gbaya {\scriptsize(gya)} & 0/1 & 0 & \multicolumn{3}{c}{Input file} \\
\raggedright Norwegian {\scriptsize(nor)} & 2/2 & 54,640 & 53.4 & 73.2 & 73.2 \\
\raggedright Occitan {\scriptsize(oci)} & 1/1 & 9,976 & 64.6 & 73.6 & 72.6 \\
\raggedright Odia {\scriptsize(ori)} & 0/1 & 0 & \multicolumn{3}{c}{Input file} \\
\raggedright Old Church Slavonic {\scriptsize(chu)} & 1/1 & 20,075 & 22.1 & 23.2 & 16.7 \\
\raggedright Old East Slavic {\scriptsize(orv)} & 4/4 & 75,765 & 35.6 & 50.3 & 53.6 \\
\raggedright Old French {\scriptsize(fro)} & 2/2 & 26,674 & 41.9 & 41.8 & 20.5 \\
\raggedright Old Georgian {\scriptsize(oge)} & 3/3 & 31,185 & 50.0 & 69.5 & 79.7 \\
\raggedright Old Irish {\scriptsize(sga)} & 2/2 & 824 & 42.8 & 44.1 & 33.0 \\
\raggedright Old Occitan {\scriptsize(pro)} & 1/2 & 9,348 & 45.0 & 48.2 & 50.5 \\
\raggedright Ottoman Turkish {\scriptsize(ota)} & 3/3 & 17,783 & 49.3 & 68.4 & 78.3 \\
\raggedright Pashto {\scriptsize(pst)} & 1/2 & 1,118 & 69.8 & 73.4 & 75.4 \\
\raggedright Persian {\scriptsize(pes)} & 2/2 & 40,127 & 80.4 & 96.6 & 97.0 \\
\raggedright Pesh {\scriptsize(pay)} & 0/1 & 0 & \multicolumn{3}{c}{Input file} \\
\raggedright Phrygian {\scriptsize(xpg)} & 0/1 & 0 & \multicolumn{3}{c}{Input file} \\
\raggedright Polish {\scriptsize(pol)} & 4/4 & 71,567 & 50.3 & 67.5 & 64.7 \\
\raggedright Pomak {\scriptsize(qpm)} & 1/1 & 8,797 & 42.9 & 85.2 & 88.2 \\
\raggedright Portuguese {\scriptsize(por)} & 10/10 & 274,394 & 70.1 & 73.7 & 47.9 \\
\raggedright Punjabi {\scriptsize(pan)} & 3/3 & 997 & 65.2 & 73.1 & 75.0 \\
\raggedright Romanian {\scriptsize(ron)} & 6/6 & 57,002 & 51.2 & 59.9 & 58.2 \\
\raggedright Ruruuli-Lunyala {\scriptsize(ruc)} & 1/1 & 8,082 & 41.4 & 52.7 & 52.7 \\
\raggedright Russian {\scriptsize(rus)} & 5/5 & 213,826 & 53.8 & 89.4 & 89.8 \\
\raggedright Sanskrit {\scriptsize(san)} & 2/2 & 21,217 & 26.6 & 80.4 & 86.7 \\
\raggedright Scottish Gaelic {\scriptsize(gla)} & 1/1 & 10,552 & 64.2 & 92.2 & 90.4 \\
\raggedright Serbian {\scriptsize(srp)} & 2/2 & 21,572 & 48.4 & 81.4 & 78.9 \\
\raggedright Shanghainese (Wu Chinese) {\scriptsize(wuu)} & 1/1 & 874 & 100.0 & 92.8 & 92.8 \\
\raggedright Sicilian {\scriptsize(scn)} & 0/1 & 0 & \multicolumn{3}{c}{Input file} \\
\raggedright Sindhi {\scriptsize(snd)} & 1/1 & 10,957 & 76.4 & 96.1 & 96.1 \\
\raggedright Sinhala {\scriptsize(sin)} & 1/2 & 668 & 56.0 & 56.0 & 51.6 \\
\raggedright Skolt Sami {\scriptsize(sms)} & 1/1 & 2,954 & 55.1 & 61.3 & 54.3 \\
\raggedright Slovak {\scriptsize(slk)} & 1/1 & 12,704 & 44.6 & 70.3 & 89.4 \\
\raggedright Slovenian {\scriptsize(slv)} & 2/2 & 36,821 & 54.1 & 89.1 & 93.3 \\
\raggedright South Levantine Arabic {\scriptsize(ajp)} & 0/1 & 0 & \multicolumn{3}{c}{Input file} \\
\raggedright Spanish {\scriptsize(spa)} & 4/4 & 96,932 & 66.9 & 78.6 & 73.1 \\
\raggedright Spanish Sign Language {\scriptsize(ssp)} & 0/1 & 0 & \multicolumn{3}{c}{Input file} \\
\raggedright Swedish {\scriptsize(swe)} & 5/6 & 69,960 & 62.3 & 72.1 & 66.1 \\
\raggedright Tagalog {\scriptsize(tgl)} & 1/2 & 1,097 & 69.6 & 74.5 & 74.3 \\
\raggedright Tamil {\scriptsize(tam)} & 1/2 & 1,989 & 48.7 & 52.4 & 58.5 \\
\raggedright Tatar {\scriptsize(tat)} & 0/1 & 0 & \multicolumn{3}{c}{Input file} \\
\raggedright Teko {\scriptsize(eme)} & 0/1 & 0 & \multicolumn{3}{c}{Input file} \\
\raggedright Telugu {\scriptsize(tel)} & 0/1 & 0 & \multicolumn{3}{c}{No gold} \\
\raggedright Thai {\scriptsize(tha)} & 1/3 & 5,739 & 99.7 & 99.7 & 99.7 \\
\raggedright Tupinamba {\scriptsize(tpn)} & 0/1 & 0 & \multicolumn{3}{c}{Input file} \\
\raggedright Turkish {\scriptsize(tur)} & 10/10 & 99,232 & 48.7 & 53.2 & 51.0 \\
\raggedright Turkish English {\scriptsize(qti)} & 0/1 & 0 & \multicolumn{3}{c}{Input file} \\
\raggedright Turkish German {\scriptsize(qtd)} & 1/1 & 14,088 & 64.9 & 85.3 & 87.0 \\
\raggedright Turkmen {\scriptsize(tuk)} & 0/0 & 0 & \multicolumn{3}{c}{No test} \\
\raggedright Ukrainian {\scriptsize(ukr)} & 3/3 & 37,049 & 54.4 & 64.5 & 65.0 \\
\raggedright Umbrian {\scriptsize(xum)} & 0/1 & 0 & \multicolumn{3}{c}{Input file} \\
\raggedright Upper Sorbian {\scriptsize(hsb)} & 0/1 & 0 & \multicolumn{3}{c}{Input file} \\
\raggedright Urdu {\scriptsize(urd)} & 1/1 & 14,537 & 75.0 & 76.9 & 75.6 \\
\raggedright Uyghur {\scriptsize(uig)} & 1/1 & 7,573 & 64.1 & 86.5 & 93.5 \\
\raggedright Uzbek {\scriptsize(uzb)} & 3/3 & 8,995 & 52.5 & 52.7 & 55.8 \\
\raggedright Veps (Vepsian) {\scriptsize(vep)} & 1/1 & 1,303 & 45.0 & 74.8 & 73.4 \\
\raggedright Vietnamese {\scriptsize(vie)} & 2/2 & 13,526 & 93.2 & 96.9 & 96.8 \\
\raggedright Warlpiri {\scriptsize(wbp)} & 0/1 & 0 & \multicolumn{3}{c}{Input file} \\
\raggedright Welsh {\scriptsize(cym)} & 1/1 & 17,026 & 61.0 & 91.4 & 89.8 \\
\raggedright Western Armenian {\scriptsize(hyw)} & 1/1 & 13,493 & 56.2 & 84.9 & 91.7 \\
\raggedright Western Sierra Puebla Nahuatl {\scriptsize(nhi)} & 0/1 & 0 & \multicolumn{3}{c}{Input file} \\
\raggedright Wolof {\scriptsize(wol)} & 1/1 & 10,403 & 73.2 & 92.1 & 91.7 \\
\raggedright Xavante {\scriptsize(xav)} & 0/1 & 0 & \multicolumn{3}{c}{Input file} \\
\raggedright Xibe {\scriptsize(sjo)} & 1/1 & 15,395 & 79.8 & 79.8 & 79.8 \\
\raggedright Yakut {\scriptsize(sah)} & 0/1 & 0 & \multicolumn{3}{c}{Input file} \\
\raggedright Yiddish {\scriptsize(yid)} & 1/1 & 9,955 & 67.6 & 91.0 & 92.1 \\
\raggedright Yoruba {\scriptsize(yor)} & 0/1 & 0 & \multicolumn{3}{c}{Input file} \\
\raggedright Yupik {\scriptsize(ess)} & 0/1 & 0 & \multicolumn{3}{c}{Input file} \\
\raggedright Zaar {\scriptsize(say)} & 1/1 & 9,965 & 61.5 & 80.0 & 80.7 \\
\raggedright Zazaki {\scriptsize(zza)} & 1/1 & 249 & 60.2 & 79.1 & 79.1 \\
\end{supertabular*}
\endgroup


\end{document}
